To isolate the specific contributions of schedule conditioning, state representation design, and gradient coupling, we conduct systematic ablation studies (Tables~\ref{tab:main}, \ref{tab:pred}, and~\ref{tab:eta}).

\paragraph{Impact of Schedule Conditioning (Ablation A2)} 
Removing the transit schedule features from the predictor and the policy input degrades forecasting clearly and control less clearly. The forecast RMSE rises from $\mAOneRmse$ to $\mATwoRmse$ and the surge-window RMSE from $\mPropSurgeRmse$ to $\mATwoSurgeRmse$ (Table~\ref{tab:pred}). For control, the like-for-like comparison is with the latent-embedding configuration with $\eta=0.1$, which also uses the embedding: its reward falls from $\mFullReward$ to $\mATwoReward$ (about $9\%$) and its latency rises from $\mFullLat$~ms to $\mATwoLat$~ms ($7.8\%$), but with three seeds the reward difference is within seed variation (its 95\% interval includes zero). Against the proposed configuration the gap is larger (reward $\mPropReward$ to $\mATwoReward$, latency $\mPropLat$~ms to $\mATwoLat$~ms), but that comparison also changes the state representation. We read these results as evidence that the timetable carries information that the historical counters do not, for forecasting; for control the evidence is suggestive only.

\paragraph{State Representation Design: Direct Schedule Features vs.\ Latent Embeddings} 
We investigate how predictive schedule knowledge is best communicated to the DRL policy: via compact, explicit schedule features $\mathbf{S}_i^{t+1}\in\mathbb{R}^8$ versus the full high-dimensional spatio-temporal latent embedding $\mathbf{Z}_{\mathrm{fused}}\in\mathbb{R}^{768}$. The results suggest the following for our setting:
\begin{itemize}
  \item \textbf{Compact Direct Representation (Proposed STAG-FedSAC)}: Directly injecting the 8 normalized schedule features (event times, passenger counts, and surge indicators) into the SAC state provides a compact, stationary, and semantically invariant representation. Because these features depend only on the deterministic timetable, transitions stored in the off-policy replay buffer $\mathcal{B}_i$ remain valid across subsequent training iterations, which should favour stable Q-function estimation; this configuration gives the best control performance among the SAC variants ($\mPropReward$, latency $\mPropLat$~ms).
  \item \textbf{Latent Embedding Injection}: In contrast, passing the uncompressed 768-dimensional latent embedding $\mathbf{Z}_{\mathrm{fused}}$ into the policy state exposes off-policy training to \emph{representation staleness}. As the ST-GCAT predictor updates its parameters $\phi$ during training, historical transitions stored in the replay buffer contain stale representations generated by earlier model checkpoints. A plausible consequence is non-stationarity in the critic's value estimates; we did not test this directly (for instance by recomputing the embeddings of stored transitions). The observed reward is lower than with the compact representation: $\mAOneReward$ with $\eta=0$ (latency $\mAOneLat$~ms) and $\mFullReward$ with $\eta=0.1$ (latency $\mFullLat$~ms), against $\mPropReward$ ($\mPropLat$~ms).
\end{itemize}
In our experiments, compact schedule features worked better than concatenating the latent embedding; we did not test other ways of using the embedding.

\paragraph{Predictor--Policy Gradient Coupling ($\eta$ Sweep)} 
We examine the effect of the bidirectional joint gradient coupling weight $\eta$ in Eq.~\eqref{eq:joint}, which backpropagates the actor objective into the ST-GCAT predictor. Table~\ref{tab:eta} evaluates $\eta \in \{0, 0.01, 0.10\}$. Among the three values tested, the decoupled objective ($\eta=0$) gives the highest reward ($\mAOneReward$) and the lowest forecast error (RMSE $\mAOneRmse$, Surge-RMSE $\mAOneSurgeRmse$). Increasing the joint gradient weight to $\eta=0.01$ and $\eta=0.10$ lowers the control reward (to $\mEtaSmallFullReward$ and $\mFullReward$) and slightly worsens the forecast accuracy (RMSE $\mEtaSmallFullRmse$ and $\mFullRmse$). 

A plausible explanation, which we did not test, is that in off-policy learning the critic evaluates historical buffer samples while the actor objective changes the predictor, so that backpropagating the policy gradient into ST-GCAT perturbs the embeddings that the stored transitions were computed with. We could not show a benefit of the joint gradient at any of the three weights.

\begin{table}[htbp]
\centering
\caption{Sensitivity to the joint gradient coupling weight $\eta$ (mean $\pm$ standard deviation over seeds; $\eta=0$ corresponds to decoupled training).}
\label{tab:eta}
\small
\begin{tabular}{llcccc}
\toprule
$\eta$ & Seeds & Reward & Lat. (ms) & RMSE & Surge-RMSE \\
\midrule
0 & 3 & 0.176{\tiny$\pm$0.009} & 118.4{\tiny$\pm$9.8} & 0.121{\tiny$\pm$0.000} & 0.163{\tiny$\pm$0.009} \\
0.1 & 3 & 0.141{\tiny$\pm$0.008} & 170.6{\tiny$\pm$8.6} & 0.127{\tiny$\pm$0.000} & 0.169{\tiny$\pm$0.006} \\
\bottomrule
\end{tabular}

\end{table}
